\chapter{Construction du système de scoring par apprentissage automatique}

\section{Construction de la variable cible}

Les données disponibles ne comportent pas d'historique réel d'incidents de paiement permettant
de définir directement une variable de défaut. Face à cette contrainte, deux options ont été
envisagées : (i) fixer un seuil arbitraire sur le score MNS2 réel (\texttt{score\_final}), ou
(ii) laisser les données elles-mêmes déterminer la coupure séparant les dossiers « sains » des
dossiers « à risque ».

La première option a été testée et écartée : sur l'échantillon réel disponible,
\texttt{score\_final} varie entre 53 et 78. Un seuil fixé a priori à 40 (valeur usuelle dans la
littérature bancaire pour des scores sur 100) ne produit \emph{aucun} dossier en défaut, rendant
tout apprentissage supervisé impossible (variance nulle de la cible).

\begin{definition}{Seuil de défaut dynamique}
Le seuil de défaut est recalculé automatiquement à chaque exécution du pipeline par un
algorithme de partitionnement non supervisé (K-means, $k=2$) appliqué à la distribution des
valeurs de \texttt{score\_final} disponibles. Soit $c_1 < c_2$ les deux centres de cluster
obtenus ; le seuil retenu est le point médian $\tau = (c_1 + c_2)/2$. Un dossier est étiqueté
« en défaut » si son score est strictement inférieur à $\tau$.
\end{definition}

Sur l'échantillon actuel (7 dossiers avec \texttt{score\_final} renseigné), cette procédure
aboutit à $\tau \approx 65{,}6$, séparant l'échantillon en deux groupes de taille comparable
(Tableau~\ref{tab:seuil-defaut}).

\begin{limite}[Cible proxy, non un défaut réel]
La variable cible ainsi construite est un \emph{proxy statistique} — une coupure relative dans
la distribution des scores disponibles — et non un historique réel de défaut de paiement. Son
interprétation doit rester prudente : elle indique quels dossiers se distinguent
\emph{relativement} aux autres dossiers de l'échantillon, non une probabilité de défaut au sens
strict (Bâlois, cf. Chapitre~2, Section 2.1). Ce point devra être discuté explicitement lors de
la soutenance et, si possible, validé ou remplacé par un historique réel avec le maître de stage.
\end{limite}

L'avantage de cette approche dynamique, par rapport à un seuil fixe, est qu'elle reste valide à
mesure que de nouveaux dossiers seront intégrés via le workflow applicatif (Chapitre~6) : le
seuil se recalcule alors automatiquement sur la distribution actualisée, sans intervention
manuelle.

\section{Variables explicatives et protocole d'évaluation}

Sur les 339 variables canoniques (Chapitre~3), 34 ont été retenues comme candidates explicatives
sur la base de leur couverture (présentes pour la majorité des dossiers) et de leur pertinence
financière (postes de bilan, de résultat, ratios de structure). Les valeurs manquantes sont
imputées par la médiane calculée sur l'échantillon d'entraînement (jamais sur l'échantillon de
test, pour éviter toute fuite d'information), et les variables sont standardisées.

\begin{definition}{Validation croisée Leave-One-Out (LOOCV)}
Pour un échantillon de $n$ observations, la validation croisée \emph{leave-one-out} consiste à
entraîner le modèle $n$ fois, chaque itération excluant une observation différente de
l'ensemble d'entraînement pour l'utiliser comme unique observation de test. C'est le seul
protocole de validation qui exploite la totalité de l'échantillon disponible tout en garantissant
qu'aucune observation testée n'a participé à l'entraînement du modèle qui la prédit — propriété
indispensable lorsque $n$ est trop restreint pour un découpage classique
entraînement/validation/test.
\end{definition}

Deux modèles sont entraînés en parallèle sur les 7 dossiers labellisés : une \textbf{forêt
aléatoire} (300 arbres, profondeur maximale 4, pondération des classes équilibrée) et une
\textbf{régression logistique} — cette dernière servant de modèle de référence pour mesurer
l'apport réel du changement de paradigme algorithmique (cf. justification théorique,
Chapitre~2, Section 2.2).

\section{Résultats}

Le Tableau~\ref{tab:proba-defaut} présente, pour chaque dossier labellisé, la probabilité de
défaut estimée par chacun des deux modèles en LOOCV.

\begin{table}[h]
\centering
\small
\begin{tabular}{@{}lccc@{}}
\toprule
\textbf{Dossier} & \textbf{P(défaut) — RF} & \textbf{P(défaut) — LogReg} & \textbf{Label réel} \\
\midrule
Feuil8  & 79.3\,\% & 98.6\,\% & Sain \\
Feuil11 & 75.0\,\% & 96.6\,\% & Sain \\
Feuil13 & 60.3\,\% & 83.0\,\% & Défaut \\
Feuil2  & 43.0\,\% & 61.1\,\% & Défaut \\
Feuil5  & 49.7\,\% & 37.0\,\% & Défaut \\
Feuil1  & 40.3\,\% & 42.7\,\% & Défaut \\
Feuil3  & 71.0\,\% & 11.1\,\% & Sain \\
\bottomrule
\end{tabular}
\caption{Probabilités de défaut estimées en LOOCV (n=7)}
\label{tab:proba-defaut}
\end{table}
\label{tab:seuil-defaut}

Plusieurs observations se dégagent. D'abord, les deux modèles se trompent nettement sur Feuil8
et Feuil11 (labellisés sains mais prédits en défaut avec forte confiance), et divergent fortement
entre eux sur Feuil3 (71\,\% pour la forêt aléatoire contre 11\,\% pour la régression logistique).
Ce désaccord inter-modèles, combiné à des erreurs de classification sur un échantillon aussi
restreint, doit être interprété comme un signal de \textbf{sur-apprentissage} plausible plutôt
que comme une réelle difficulté de prédiction : avec 6 observations d'entraînement à chaque
itération LOOCV, la capacité de généralisation des deux modèles est structurellement limitée.

\begin{limite}[Performance non généralisable]
Avec $n=7$ dossiers labellisés, les métriques de performance (taux de bonne classification, AUC)
calculées sur ce LOOCV ne doivent en aucun cas être présentées comme une estimation fiable de la
performance du système sur de nouveaux dossiers. Elles sont rapportées à titre indicatif, pour
comparer les deux modèles entre eux sur cet échantillon précis, conformément au cadre posé au
Chapitre~2 (Section 2.3).
\end{limite}

\section{Importance des variables}

Le Tableau~\ref{tab:feature-importance} présente les dix variables les plus déterminantes selon
la diminution moyenne d'impureté (MDI) du modèle de forêt aléatoire entraîné sur l'ensemble de
l'échantillon labellisé.

\begin{table}[h]
\centering
\small
\begin{tabular}{@{}lc@{}}
\toprule
\textbf{Variable} & \textbf{Importance (MDI)} \\
\midrule
Matériel et agencement & 0.076 \\
Banque, caisse, CCP & 0.070 \\
Financement permanent & 0.066 \\
Rentabilité globale & 0.049 \\
Capacité de remboursement & 0.048 \\
Actif immobilisé net & 0.048 \\
Capital social personnel & 0.044 \\
Trésorerie nette & 0.043 \\
Consommation de l'exercice & 0.040 \\
Réserves & 0.038 \\
\bottomrule
\end{tabular}
\caption{Facteurs les plus déterminants (importance MDI, Random Forest)}
\label{tab:feature-importance}
\end{table}

La prédominance de variables liées à la \textbf{structure de l'actif immobilisé} (matériel et
agencement, actif immobilisé net) et à la \textbf{liquidité immédiate} (banque/caisse/CCP,
trésorerie nette) est cohérente avec l'intuition financière : ce sont des indicateurs directs de
la capacité de l'entreprise à honorer ses engagements à court terme. La présence du
\emph{financement permanent} en troisième position souligne également le rôle de la structure de
financement à long terme dans la discrimination des dossiers.

\begin{limite}[Importance par permutation non informative]
L'importance par permutation, calculée en complément de l'importance MDI, retourne des valeurs
quasi nulles pour l'ensemble des variables sur cet échantillon. Ce résultat s'explique par la
taille de l'échantillon labellisé (7 observations) : la permutation d'une variable sur un jeu de
test aussi restreint ne modifie pas suffisamment la performance mesurée pour être détectée de
façon fiable. Seule l'importance MDI est donc interprétée dans ce chapitre.
\end{limite}

\section{Synthèse}

Le système de scoring par apprentissage automatique met en évidence des facteurs financiers
cohérents avec la théorie du risque de crédit, mais ses prédictions individuelles doivent être
interprétées avec une grande prudence compte tenu de la taille de l'échantillon. Le
Chapitre~6 combine ce score avec le score expert (Chapitre~4) au sein d'un moteur de décision
unique, et confronte les résultats obtenus à la notation MNS2 réelle.
